\chapter{Granice i ciągłość funkcji wielu zmiennych}
\label{ma2:ch:multivariable-continuity}

Ogólne wyniki o granicach i ciągłości w przestrzeniach
metrycznych znamy z AM I, rozdziały 14--16. Przytaczamy je
z odsyłaczami do dowodów; skupimy się na konsekwencjach
wielowymiarowych, przykładach i kontrprzykładach.

\section{Przestrzenie współrzędnych i granice}

Niech $n,m\ge1$, $A\subseteq\R^n$, $f:A\to\R^m$,
$f=(f_1,\ldots,f_m)$, gdzie $f_j:A\to\R$.
Domyślnie stosujemy normę euklidesową
$\|x\|_2=(\sum_{j=1}^n x_j^2)^{1/2}$. Z twierdzeń
\ref{ma2:thm:equiv} i \ref{ma2:thm:finite}
wynika niezależność zbieżności, ciągłości i zupełności
od wyboru normy. W szczególności
\[
 \|x\|_\infty\le\|x\|_2\le\sqrt n\,\|x\|_\infty,
 \qquad \|x\|_\infty=\max_{1\le j\le n}|x_j|.
\]
Pierwsza nierówność wynika z uwzględnienia największego
składnika sumy kwadratów, druga z ograniczenia każdego
składnika przez $\|x\|_\infty^2$.

Przypomnijmy definicje AM I, 14.1.1 i 14.1.3.
Punkt $p\in\R^n$ jest punktem skupienia zbioru $A$, gdy
dla każdego $\delta>0$ istnieje $x\in A$ z
$0<\|x-p\|_2<\delta$. Zbiór tych punktów oznaczamy $A'$.
Dla $p\in A'$ i $L\in\R^m$ warunek
$\lim_{A\ni x\to p}f(x)=L$ oznacza
\[
 \forall\varepsilon>0\ \exists\delta>0\ \forall x\in A:
 \quad 0<\|x-p\|_2<\delta\Longrightarrow
 \|f(x)-L\|_2<\varepsilon.
\]
Punkt $p$ nie musi należeć do $A$.

\begin{twierdzenie}[Kryteria granicy]
\label{ma2:thm:multi-limit}
Niech $p\in A'$ i $L=(L_1,\ldots,L_m)\in\R^m$.
Następujące warunki są równoważne:
\begin{enumerate}[label=(\roman*)]
 \item $\lim_{A\ni x\to p}f(x)=L$;
 \item dla każdego ciągu $(x_k)$ w $A\setminus\{p\}$
 zbieżnego do $p$ zachodzi $f(x_k)\to L$;
 \item dla każdego $j=1,\ldots,m$ zachodzi
 $\lim_{A\ni x\to p}f_j(x)=L_j$.
\end{enumerate}
\end{twierdzenie}
\begin{proof}
Równoważność (i) i (ii) jest szczególnym przypadkiem
twierdzenia AM I, 14.1.5.
Dla (i)$\Rightarrow$(iii) korzystamy z nierówności
$|f_j(x)-L_j|\le\|f(x)-L\|_2$.
Dla (iii)$\Rightarrow$(i) ustalamy $\varepsilon>0$ i
dla każdej współrzędnej wybieramy $\delta_j>0$ dla
dokładności $\varepsilon/\sqrt m$.
Przy $\delta=\min_j\delta_j>0$ i
$0<\|x-p\|_2<\delta$ mamy
\[
 \|f(x)-L\|_2^2=\sum_{j=1}^m|f_j(x)-L_j|^2
 < m(\varepsilon/\sqrt m)^2=\varepsilon^2.
\]
To daje (i).
\end{proof}

\begin{wniosek}
\label{ma2:cor:multi-limit-cauchy}
Dla $p\in A'$ granica $f$ w $p$ istnieje w $\R^m$ wtedy
i tylko wtedy, gdy dla każdego $\varepsilon>0$ istnieje
$\delta>0$ takie, że $\|f(x)-f(y)\|_2<\varepsilon$
dla wszystkich $x,y\in A$ spełniających
$0<\|x-p\|_2<\delta$, $0<\|y-p\|_2<\delta$.
\end{wniosek}
Jest to kryterium Cauchy'ego z AM I, twierdzenie 14.2.4:
jego założenie zupełności przestrzeni wartości jest spełnione
na mocy twierdzenia \ref{ma2:thm:finite}.
Rachunek granic skalarnych i twierdzenie o trzech funkcjach
stosujemy z AM I, twierdzenia 14.3.1--2, a do odwzorowań
wektorowych stosujemy kryterium współrzędnych.

\section{Granice wzdłuż krzywych i granice iterowane}

Jeżeli $\lim_{A\ni x\to p}f(x)=L$ oraz
$\gamma:(0,r)\to A\setminus\{p\}$ spełnia
$\lim_{t\to0^+}\gamma(t)=p$, to
$\lim_{t\to0^+}f(\gamma(t))=L$ na mocy twierdzenia
AM I, 14.2.2: argument $\gamma(t)$ unika $p$.
Różne granice wzdłuż dwóch takich krzywych wykluczają
granicę wspólną.

\begin{przyklad}[Proste nie wystarczają]
\label{ma2:ex:multivariable-paths}
Niech $f:\R^2\to\R$ będzie dana wzorem
\[
 f(x,y)=
 \begin{cases}
 \dfrac{x^2y}{x^4+y^2},&(x,y)\ne(0,0),\\
 0,&(x,y)=(0,0).
 \end{cases}
\]
Dla $(a,b)\ne(0,0)$, $t\ne0$ i $b\ne0$ mamy
\[
 f(ta,tb)=\frac{ta^2b}{t^2a^4+b^2}\longrightarrow0.
\]
Gdy $b=0$, wartości na prostej są zerowe.
Ograniczenie do każdej prostej przez początek jest
ciągłe w początku. Jednak dla $t\ne0$ mamy
$f(t,t^2)=1/2$, podczas gdy $f(t,0)=0$.
Granica w $(0,0)$ nie istnieje, a $f$ nie jest tam ciągła.
\end{przyklad}

\begin{figure}[htbp]
\centering
\begin{tikzpicture}[>=Stealth,scale=1.25]
 \draw[->] (-1.6,0)--(1.7,0) node[right] {$x$};
 \draw[->] (0,-.5)--(0,1.9) node[above] {$y$};
 \draw[blue!70!black,thick,domain=-1.2:1.2,samples=65]
 plot (\x,{\x*\x});
 \draw[red!70!black,very thick,->] (1.25,0)--(.45,0);
 \draw[blue!70!black,very thick,->]
 plot[domain=1.05:.65,samples=20] (\x,{\x*\x});
 \fill (0,0) circle (1.4pt) node[below left] {$0$};
 \node[blue!70!black,anchor=west] at (1.0,1.52) {$y=x^2$};
 \node[red!70!black] at (.9,-.3) {$y=0$};
 \node[blue!70!black,anchor=east] at (-.95,1.5) {$f=1/2$};
 \node[red!70!black] at (-.9,.2) {$f=0$};
\end{tikzpicture}
\caption{Dwie drogi do początku układu z różnymi granicami
funkcji z przykładu \ref{ma2:ex:multivariable-paths}.
Wartości na paraboli dotyczą punktów różnych od początku.}
\label{ma2:fig:multivariable-paths}
\end{figure}

\begin{przyklad}[Granice iterowane]
Dla tej samej funkcji i każdego ustalonego $x$ zachodzi
$\lim_{y\to0}f(x,y)=0$: dla $x\ne0$ mianownik dąży
do $x^4>0$, a dla $x=0$ funkcja jest zerowa.
Dla każdego ustalonego $y$ analogicznie
$\lim_{x\to0}f(x,y)=0$ (dla $y\ne0$ mianownik dąży
do $y^2>0$, dla $y=0$ wartości są zerowe).
Obie granice iterowane są więc równe $0$,
ale granica wspólna nie istnieje.
\end{przyklad}

\begin{twierdzenie}[Jednostajność względem kierunku]
\label{ma2:thm:radial-limit}
Niech $p\in\R^n$, $\rho>0$,
$B(p,\rho)\setminus\{p\}\subseteq A$ i
$S^{n-1}=\{u\in\R^n:\|u\|_2=1\}$.
Wtedy $\lim_{A\ni x\to p}f(x)=L$ zachodzi wtedy
i tylko wtedy, gdy
\[
 \lim_{r\to0^+}\sup_{u\in S^{n-1}}\|f(p+ru)-L\|_2=0,
 \qquad 0<r<\rho.
\]
Dopuszczamy supremum $+\infty$ dla niektórych $r$.
\end{twierdzenie}
\begin{proof}
Jeśli granica istnieje, dla $\varepsilon>0$ wybieramy
$\delta>0$ z definicji granicy dla dokładności $\varepsilon/2$.
Dla $0<r<\min(\rho,\delta)$ i każdego $u\in S^{n-1}$
mamy $\|p+ru-p\|_2=r<\delta$, zatem supremum nie przekracza
$\varepsilon/2$. To dowodzi granicy supremum.
Odwrotnie, dla $\varepsilon>0$ wybieramy $\delta\in(0,\rho)$,
dla którego supremum jest mniejsze od $\varepsilon$ dla
$0<r<\delta$. Dla $x\in A$ z $0<\|x-p\|_2<\delta$
kładziemy $r=\|x-p\|_2$, $u=(x-p)/r$.
Wówczas $u\in S^{n-1}$, $x=p+ru$ i
$\|f(x)-L\|_2<\varepsilon$.
\end{proof}

\begin{przyklad}
Dla $(x,y)\ne(0,0)$ niech $g(x,y)=x^2y/(x^2+y^2)$.
Ponieważ $x^2/(x^2+y^2)\le1$, mamy
$|g(x,y)|\le |y|\le\|(x,y)\|_2$.
Stąd granica jest równa $0$; przedłużenie $g(0,0)=0$
jest ciągłe. Oszacowanie nie zależy od kierunku,
w przeciwieństwie do przykładu \ref{ma2:ex:multivariable-paths}.
\end{przyklad}

\section{Ciągłość i działania na odwzorowaniach}

Dla $p\in A$ ciągłość oznacza
\[
 \forall\varepsilon>0\ \exists\delta>0\ \forall x\in A:
 \quad\|x-p\|_2<\delta\Longrightarrow
 \|f(x)-f(p)\|_2<\varepsilon.
\]
Jest to definicja AM I, 15.1.1. W punkcie skupienia
jest równoważna istnieniu granicy równej $f(p)$,
a w punkcie izolowanym zachodzi automatycznie.
Twierdzenie AM I, 15.1.3, daje równoważność z warunkiem
$f(x_k)\to f(p)$ dla każdego ciągu $(x_k)$ w $A$
zbieżnego do $p$. Z twierdzenia \ref{ma2:thm:multi-limit}
i przypadku punktów izolowanych wynika, że ciągłość
$f$ jest równoważna ciągłości wszystkich współrzędnych.

Według twierdzenia AM I, 15.2.4, $f$ jest ciągłe na $A$
wtedy i tylko wtedy, gdy przeciwobrazy zbiorów otwartych
w $\R^m$ są otwarte względnie w $A$; równoważnie
przeciwobrazy zbiorów domkniętych w $\R^m$ są domknięte
względnie w $A$. Jeśli $f:A\to B$ jest ciągłe w $p\in A$,
$B\subseteq\R^m$, a $g:B\to\R^k$ jest ciągłe w $f(p)$,
to $g\circ f$ jest ciągłe w $p$ (AM I, 15.3.1).

\begin{wniosek}
\label{ma2:cor:multi-polynomials}
Wielomiany $P:\R^n\to\R$ są ciągłe, a ilorazy
$P/Q$ wielomianów są ciągłe na otwartym zbiorze
$\{x:Q(x)\ne0\}$.
\end{wniosek}
\begin{proof}
Rzuty współrzędnych są ciągłe, gdyż
$|x_j-y_j|\le\|x-y\|_2$. Stałe są ciągłe.
Każdy wielomian jest skończoną sumą iloczynów stałych
i rzutów, więc stosujemy rachunek funkcji ciągłych
(AM I, twierdzenie 15.3.2).
Zbiór, gdzie $Q\ne0$, jest przeciwobrazem otwartego
zbioru $\R\setminus\{0\}$, a ciągłość ilorazu na nim
wynika z tego samego twierdzenia.
\end{proof}

\begin{przyklad}
Odwzorowanie $F:\R^2\to\R^3$,
$F(x,y)=(x+y,xy,\sin(x^2+y^2))$, jest ciągłe:
pierwsze dwie współrzędne są wielomianami, trzecia
jest złożeniem ciągłego sinusa z wielomianem
(AM I, rozdział 18).
\end{przyklad}

\section{Zwartość i jednostajna ciągłość}

\begin{twierdzenie}
\label{ma2:thm:euclidean-compact}
Zbiór $K\subseteq\R^n$ jest ciągowo zwarty wtedy i tylko wtedy,
gdy jest domknięty i ograniczony.
\end{twierdzenie}
\begin{proof}
Jeśli $K$ jest ciągowo zwarty, to każdy ciąg w $K$
zbieżny w $\R^n$ ma podciąg zbieżny do punktu $K$.
Jednoznaczność granicy daje domkniętość przez kryterium
ciągowe (AM I, lemat 7.5.2).
Gdyby $K$ nie był ograniczony, można by wybrać
$x_k\in K$ z $\|x_k\|_2>k$. Żaden podciąg nie byłby
zbieżny, gdyż każdy zbieżny ciąg jest ograniczony.
Otrzymujemy sprzeczność.

Niech teraz $K$ będzie domknięty i ograniczony,
a $(x_k)$ ciągiem w $K$. Stosujemy twierdzenie
Bolzana--Weierstrassa (AM I, 7.2.4) do ograniczonego
ciągu pierwszych współrzędnych, wybierając podciąg
zbieżny w tej współrzędnej. Z niego wybieramy podciąg
zbieżny w drugiej współrzędnej itd. Po $n$ krokach
otrzymujemy jeden podciąg zbieżny w każdej współrzędnej;
wybieranie dalszych podciągów nie zmienia wcześniejszych
granic. Nierówność między normą maksimum i euklidesową
daje zbieżność wektorową. Domkniętość $K$ zapewnia
przynależność granicy do $K$. Dla pustego $K$
warunek ciągowej zwartości zachodzi w sposób pusty.
\end{proof}

W przestrzeniach metrycznych zwartość i zwartość ciągowa
są równoważne (topologia ogólna, rozdział 07).
Przytaczamy trzy wyniki AM I:
\begin{enumerate}[label=(\roman*)]
 \item Jeśli $K\subseteq\R^n$ jest domknięty i ograniczony,
 a $f:K\to\R^m$ ciągłe, to $f(K)$ jest ciągowo zwarty
 (AM I, 15.4.3), zatem domknięty i ograniczony.
 \item Jeśli dodatkowo $K\ne\varnothing$ i $m=1$, to $f$
 jest ograniczona i osiąga kres dolny oraz górny
 (twierdzenie Weierstrassa, AM I, 15.4.4).
 \item Jeśli $K\subseteq\R^n$ jest domknięty i ograniczony,
 a $f:K\to\R^m$ ciągłe, to $f$ jest jednostajnie ciągłe
 (twierdzenie Heinego--Cantora, AM I, 16.2.1).
\end{enumerate}
Hipotezę ciągowej zwartości we wszystkich trzech wynikach
zapewnia twierdzenie \ref{ma2:thm:euclidean-compact}.
Jednostajna ciągłość oznacza, że dla każdego $\varepsilon>0$
istnieje $\delta>0$ wspólne dla wszystkich par $x,y\in K$
takie, że $\|x-y\|_2<\delta$ implikuje
$\|f(x)-f(y)\|_2<\varepsilon$.

\begin{przyklad}
Funkcja $P(x,y)=x^2+xy+y^2$ jest ciągła na
domkniętej kuli jednostkowej $K\subseteq\R^2$.
Osiąga tam minimum i maksimum oraz jest jednostajnie ciągła.
Na całej płaszczyźnie nie jest jednostajnie ciągła:
dla $u_k=(k,0)$, $v_k=(k+1/k,0)$ mamy
$\|u_k-v_k\|_2=1/k\to0$, lecz
$P(v_k)-P(u_k)=2+1/k^2$, co przeczy definicji
jednostajnej ciągłości z $\varepsilon=1$.
\end{przyklad}

\begin{przyklad}
Dla $x,y\in B(0,R)$, $R>0$, nierówność trójkąta daje
\[
 |\|x\|_2^2-\|y\|_2^2|
 =|\|x\|_2-\|y\|_2|(\|x\|_2+\|y\|_2)
 \le 2R\|x-y\|_2.
\]
Kwadrat normy spełnia warunek Lipschitza na każdej takiej
kuli, więc jest tam jednostajnie ciągły (AM I, 16.3.3).
\end{przyklad}

\begin{wniosek}[Przedłużenie do domknięcia]
\label{ma2:cor:multi-uniform-extension}
Jeśli $A\subseteq\R^n$ i $f:A\to\R^m$ jest
jednostajnie ciągłe, to istnieje dokładnie jedno
ciągłe przedłużenie $F:\overline A\to\R^m$.
Jest ono jednostajnie ciągłe, a jeśli dla $x,y\in A$
zachodzi $\|f(x)-f(y)\|_2\le C\|x-y\|_2^\alpha$,
$C\ge0$, $0<\alpha\le1$, to ta sama nierówność
z tymi samymi stałymi zachodzi dla $F$ na $\overline A$.
\end{wniosek}
Jest to twierdzenie AM I, 16.6.3, ponieważ $\R^m$ jest
zupełna (twierdzenie \ref{ma2:thm:finite}).
Wartość przedłużenia w $p$ jest granicą $f(x_k)$ dla
dowolnego ciągu $x_k\in A$ zbieżnego do $p$.

Przykład \ref{ma2:ex:multivariable-paths} będzie ważny
przy badaniu pochodnych kierunkowych: regularność
na każdej prostej nie gwarantuje nawet ciągłości.
